\documentclass[11pt]{article}
\usepackage[a4paper,margin=26mm]{geometry}
\usepackage{amsmath,amssymb,mathtools}
\usepackage{unicode-math}
\setmainfont{Latin Modern Roman}
\setmathfont{Latin Modern Math}
\usepackage{xurl,hyperref}
\hypersetup{hidelinks,pdftitle={UnsolvedMath Problem 10400007}}
\newcommand{\sref}[2]{\hyperlink{source:#1:#2}{[S#2]}}
\newcommand{\cataloguescope}[1]{#1}
\begin{document}
% UnsolvedMath ID 10400007; source catalogue code AMR-103-0007
\setcounter{section}{10400006}
\section{Problem 1.7 — (J.}\label{10400007}
{\small\sffamily UnsolvedMath ID 10400007 · Topology}
\subsection{Definitions and mathematical statement}

\paragraph{Shared notation.}$\mathbb N=\{1,2,\ldots\}$, and $\mathbb Z,\mathbb Q,\mathbb R,\mathbb C$ denote the integers, rationals, reals and complex numbers. Any other conventions are specified locally.


A knot or link is a smooth oriented embedding of one or more circles in the oriented $3$-sphere, considered up to orientation-preserving ambient isotopy. Write $O$ for the unknot, which bounds an embedded disk. The Jones polynomial is normalized by $V_O(t)=1$ and
\[t^{-1}V_{L_+}(t)-tV_{L_-}(t)=(t^{1/2}-t^{-1/2})V_{L_0}(t),\]
where the three oriented links agree outside a ball and inside it have respectively a positive crossing, a negative crossing, and the oriented smoothing. For knots $V_K(t)\in\mathbb Z[t,t^{-1}]$.
\paragraph{Statement recorded in the supplied source.}
A local coefficient system on a space $X$ is a locally constant family of vector spaces, equivalently a representation of $\pi_1(X)$ on a vector space; its twisted homology is computed using this representation in the chain complex. Construct an intrinsic definition of $V_K$ using homology with local coefficients. Possible spaces include suitable configuration spaces associated directly to the knot or its complement, but the construction must be specified and proved to recover the Jones polynomial without requiring a chosen braid presentation. \sref{10400007}{2}
\subsection{Short English statement}
Construct the Jones polynomial intrinsically from homology with local coefficients attached to the knot.
\subsection{Sources}
\begin{description}
\item[\hypertarget{source:10400007:1}{S1}] ULAM.AI and the UnsolvedMath Contributors, \emph{UnsolvedMath: A Curated Collection of Open Mathematics Problems}, record 10400007 (AMR-103-0007). CC BY 4.0. \url{https://huggingface.co/datasets/ulamai/UnsolvedMath}.
\item[\hypertarget{source:10400007:2}{S2}] T. Ohtsuki (editor), \emph{Problems on invariants of knots and 3-manifolds}, Geometry \& Topology Monographs 4 (2002), pp. 377--572 (published 2004), item 1.7 and its surrounding definitions and remarks. \url{https://arxiv.org/abs/math/0406190v1}.
\end{description}
\label{end:10400007}
\end{document}
